\documentclass[10pt]{article}
\providecommand{\FIGDIR}{fig}
\newcommand{\DIGESTDATE}{28 August 2026}
% ===================== PM Research Watch - shared preamble =====================
\usepackage[T1]{fontenc}
\usepackage[utf8]{inputenc}
\usepackage{amsmath,amssymb}
\usepackage{tgpagella}
\usepackage{tgheros}
\renewcommand{\sfdefault}{qhv}
\renewcommand{\ttdefault}{lmtt}
\usepackage{microtype}
\usepackage[english]{babel}

\usepackage[a4paper,top=17mm,bottom=17mm,left=17mm,right=17mm,
            headsep=6mm,footskip=10mm]{geometry}
\usepackage{graphicx}
\usepackage{xcolor}
\usepackage{colortbl}
\usepackage{tikz}
\usetikzlibrary{calc}
\usepackage[most]{tcolorbox}
\usepackage{titlesec}
\usepackage{fancyhdr}
\usepackage{booktabs}
\usepackage{array}
\usepackage{longtable}
\usepackage{enumitem}
\usepackage{multicol}
\usepackage{ragged2e}
\usepackage[hidelinks]{hyperref}

% ------------------------------------------------------------------ palette
\definecolor{Ink}    {HTML}{14262E}
\definecolor{Deep}   {HTML}{0F4C5C}
\definecolor{Teal}   {HTML}{1F7A8C}
\definecolor{Sky}    {HTML}{4FA3C4}
\definecolor{Sage}   {HTML}{6FA287}
\definecolor{Amber}  {HTML}{E3A008}
\definecolor{Clay}   {HTML}{B5643C}
\definecolor{Coral}  {HTML}{D1495B}
\definecolor{Violet} {HTML}{6A5B9E}
\definecolor{Slate}  {HTML}{4A6572}
\definecolor{Paper}  {HTML}{FBF9F5}
\definecolor{Mist}   {HTML}{EEEAE2}
\definecolor{Rule}   {HTML}{D9D3C9}

\pagecolor{Paper}
\color{Ink}

\hypersetup{colorlinks=true, linkcolor=Deep, urlcolor=Teal, citecolor=Teal,
            pdftitle={PM Research Watch}, pdfauthor={Automated daily digest}}

% ------------------------------------------------------------------ headings
\titleformat{\section}
  {\sffamily\bfseries\large\color{Deep}}{}{0pt}
  {}[{\vspace{1.2mm}\color{Amber}\titlerule[1.1pt]}]
\titlespacing*{\section}{0pt}{5.5mm}{2.8mm}

\titleformat{\subsection}
  {\sffamily\bfseries\normalsize\color{Teal}}{}{0pt}{}
\titlespacing*{\subsection}{0pt}{3.6mm}{1.4mm}

\setlist[itemize]{leftmargin=4.2mm, itemsep=1.1mm, topsep=1.3mm, parsep=0pt,
                  label=\textcolor{Teal}{\footnotesize$\blacktriangleright$}}

% ------------------------------------------------------------------ header/footer
\pagestyle{fancy}
\fancyhf{}
\renewcommand{\headrulewidth}{0pt}
\renewcommand{\footrulewidth}{0pt}
\fancyhead[L]{\sffamily\scriptsize\color{Slate}PM RESEARCH WATCH\;\textcolor{Rule}{\textbar}\;\DIGESTDATE}
\fancyhead[R]{\sffamily\scriptsize\color{Slate}Particulate matter monitoring \& health effects}
\fancyfoot[C]{\sffamily\scriptsize\color{Slate}\thepage}

% ------------------------------------------------------------------ components
% metric tile
\newcommand{\metric}[3]{%
  \begin{tcolorbox}[enhanced, colback=white, colframe=Rule, boxrule=0.5pt,
    arc=1.6mm, left=2.2mm, right=2.2mm, top=1.8mm, bottom=1.8mm,
    borderline west={1.7mm}{0pt}{#3}, nobeforeafter, valign=center,
    height=17mm, width=\linewidth]
    \RaggedRight
    {\sffamily\bfseries\fontsize{17}{17}\selectfont\color{#3}#1}\\[0.5mm]
    {\sffamily\fontsize{6.5}{7.4}\selectfont\color{Slate}#2}
  \end{tcolorbox}}

% pull-quote / key insight box
\newtcolorbox{keybox}[1][]{enhanced, colback=Mist, colframe=Mist,
  arc=1.6mm, left=3.2mm, right=3.2mm, top=2.4mm, bottom=2.4mm,
  borderline west={1.7mm}{0pt}{Amber}, #1}

% subtopic band
\newcommand{\band}[2]{%
  \vspace{3.2mm}\noindent
  \begin{tcolorbox}[enhanced, colback=#2, colframe=#2, arc=1.2mm,
    left=2.6mm, right=2.6mm, top=1.1mm, bottom=1.1mm, boxrule=0pt,
    width=\linewidth, nobeforeafter]
    {\sffamily\bfseries\small\color{white}#1}
  \end{tcolorbox}\vspace{1.6mm}}

% one paper entry: {short cite}{title}{journal}{doi}{body}
\newcommand{\paperentry}[5]{%
  \par\vspace{1.9mm}\noindent
  {\sffamily\bfseries\footnotesize\color{Ink}#1}\;%
  {\sffamily\fontsize{7}{8}\selectfont\color{Slate}#3}\\[0.4mm]
  {\footnotesize\itshape\color{Deep}#2}\\[0.6mm]
  {\footnotesize #4}\\[0.3mm]
  {\fontsize{6.6}{7.6}\selectfont\color{Slate}\texttt{doi:}\,\href{https://doi.org/#5}{\texttt{\detokenize{#5}}}}
  \par}

% caption style
\newcommand{\figcap}[1]{\par\vspace{1.0mm}
  {\fontsize{7.2}{8.6}\selectfont\color{Slate}\textsf{\textbf{Fig.}}\; #1\par}\vspace{1.5mm}}

\providecommand{\permil}{\textperthousand}

\setlength{\parindent}{0pt}
\setlength{\parskip}{0pt}
\renewcommand{\arraystretch}{1.18}


\begin{document}

% ==================================================================== MASTHEAD
\thispagestyle{empty}
\begin{tikzpicture}[remember picture, overlay]
  \fill[Deep] (current page.north west) rectangle
      ($(current page.north east)+(0,-67mm)$);
  \fill[Amber] ($(current page.north west)+(0,-67mm)$) rectangle
      ($(current page.north east)+(0,-68.4mm)$);
  \foreach \x/\y/\r in {22/12/0.9, 41/26/1.5, 63/9/0.7, 88/31/1.1, 112/17/1.9,
                        138/29/0.8, 159/13/1.3, 178/54/1.0, 196/21/1.6,
                        30/44/0.6, 74/57/1.2, 125/48/0.9, 168/6/0.8, 191/59/0.7,
                        52/18/1.0, 100/60/0.6, 148/20/1.1, 66/49/1.4, 205/45/0.8}
    \fill[Sky, opacity=0.30] ($(current page.north west)+(\x mm,-\y mm)$) circle (\r mm);
\end{tikzpicture}

\vspace*{4mm}
{\sffamily\fontsize{7.6}{9}\selectfont\color{Sky}%
 AUTOMATED DAILY LITERATURE DIGEST\;\textbullet\;PhD RESEARCH WATCH\par}
\vspace{2.2mm}
{\sffamily\bfseries\fontsize{27}{29}\selectfont\color{white}%
 Particulate Matter\par}
{\sffamily\fontsize{27}{29}\selectfont\color{Sky}%
 Monitoring \& Health Effects\par}
\vspace{5.0mm}
{\sffamily\fontsize{8.4}{10}\selectfont\color{white}%
 Issue 2026-08-28\;\;\textcolor{Amber}{\textbar}\;\;Entry window 28 August 2026\;\;%
 \textcolor{Amber}{\textbar}\;\;19 records screened in scope\par}
\vspace{18mm}

% -------------------------------------------------------------- metric strip
\noindent
\begin{minipage}[t]{0.190\linewidth}\metric{19}{RECORDS IN SCOPE\\(54 EXCLUDED)}{Deep}\end{minipage}\hfill
\begin{minipage}[t]{0.190\linewidth}\metric{8}{SUBTOPIC CLUSTERS\\OCCUPIED (OF 10)}{Teal}\end{minipage}\hfill
\begin{minipage}[t]{0.190\linewidth}\metric{8}{RECORDS WITH NO\\HUMAN ENDPOINT}{Sky}\end{minipage}\hfill
\begin{minipage}[t]{0.190\linewidth}\metric{4}{EFFECT ESTIMATES\\FROM 2 STUDIES}{Amber}\end{minipage}\hfill
\begin{minipage}[t]{0.190\linewidth}\metric{6}{TIER-A\\RECORDS}{Clay}\end{minipage}

\vspace{6mm}

% -------------------------------------------------------------- signal of day
\section*{\sffamily\bfseries\color{Deep}Signal of the day}
\vspace{-1.5mm}{\color{Amber}\rule{\linewidth}{1.1pt}}\vspace{2.5mm}

\begin{keybox}
{\sffamily\bfseries\small\color{Deep}Sulfate carries three-quarters of the PM$_{2.5}$
mixture effect on infarct mortality --- and, separately, coating chemistry moves black
carbon's mass absorption cross section by $\sim$1\,m$^2$/g overnight. Both results
penalise mass-only instruments.}\\[1.4mm]
\footnotesize
Li et al.\ (\textit{Environ Sci Technol Lett}) fit a causal mixture model to county-level
US acute myocardial infarction mortality against PM$_{2.5}$ and its EC / OC / SO$_4^{2-}$
/ NO$_3^-$ constituents in a difference-in-differences design. Total mass gives
\textbf{RR 1.019 (95\% CI 1.013--1.025) per IQR}; the constituent mixture gives
\textbf{RR 1.116 (1.101--1.130) per quartile}, with \textbf{SO$_4^{2-}$ weighted 0.762} of
that total and an east--west split (SO$_4^{2-}$ 0.707 in the east, OC 0.485 in the west).
The mass estimate is not wrong, it is simply a diluted average over a mixture whose
components are not interchangeable. Zhang et al.\ (\textit{ACP}) supply the
instrumentation-side counterpart from Shenzhen: SP2 + SPAMS single-particle measurement
shows nighttime heterogeneous and daytime photochemical aging proceed by different
chemistry (nitrate-first vs.\ nitrate-then-sulfate) yet converge on nearly the same
\textbf{MAC enhancement at 532\,nm --- $\sim$0.8\,m$^2$/g overnight, $\sim$1.0\,m$^2$/g
by day} on a 1.2\,$\mu$g/m$^3$ rBC background with 155--170\,nm core MMD. For a
distributed low-cost network this is the operative number: an optical channel calibrated
against a fixed mass-scattering assumption absorbs the whole diurnal aging cycle as
apparent concentration drift.
\end{keybox}

\vspace{3mm}

% -------------------------------------------------------------- exec brief
\section*{\sffamily\bfseries\color{Deep}What changed today}
\vspace{-1.5mm}{\color{Amber}\rule{\linewidth}{1.1pt}}\vspace{2.5mm}

\footnotesize
\begin{itemize}
\item \textbf{Exposure assessment is the largest cluster (5 records), and for once it is
not all epidemiology.} It holds a satellite--surface fusion product, a new summary
statistic, a filter-composition campaign, a wave-channel source-process experiment and a
bioaerosol survey. The common thread is that all five are about \emph{how the exposure
number is constructed}, not what it does to people.

\item \textbf{A 1\,km daily PM$_{2.5}$ surface for Bangladesh now exists.} Cui et al.\
gap-fill MAIAC AOD --- more than 60\% of retrievals are missing --- with LightGBM, then
push the complete AOD through an attention-based TabNet to ground level. Held-out
$R^2=0.95$, RMSE 17.41\,$\mu$g/m$^3$ over a national mean of 68--75\,$\mu$g/m$^3$. The
RMSE is roughly a quarter of the mean; that is defensible at these concentrations but it
is not a monitoring-grade product, and the paper's own DHS linkage is ecological.

\item \textbf{Citizen-science monitoring gets its first proper critical review in this
watch.} Marcon et al.\ (\textit{Curr Environ Health Rep}) is explicit that AI-assisted
calibration and distributed sensor networks are the live methodological frontier, and
equally explicit that representativeness, not sensor accuracy, is the binding constraint
on policy uptake --- an ordering worth taking seriously.

\item \textbf{A threshold-free concentration-frequency metric.} Hopke et al.\ define the
PM$_{2.5}$ Magnitude--Rank Index (largest $k$ with $\geq k$ days at $\geq k$\,$\mu$g/m$^3$),
applied to 918 US MSA-years. It correlates 0.83--0.96 with the annual mean, the 98th
percentile and exceedance counts --- which is both the selling point and the problem: a
metric that correlates 0.96 with days $\geq15$\,$\mu$g/m$^3$ has to earn its place on
interpretability, not information content.

\item \textbf{Low-cost sensors as the primary instrument in a health-risk paper.}
Wongsasuluk et al.\ run five DustBoy stations for a year in a Chiang Mai road underpass
(385{,}622 records) and build a stepwise MLR forecast. No co-location against a reference
monitor is reported in the abstract, so the risk estimates inherit whatever bias the
DustBoy units carry in a high-humidity, high-loading tunnel --- exactly the regime where
nephelometric error is largest.

\item \textbf{Policy-scale co-benefit modelling.} Zhang et al.\ (\textit{Nat Commun})
put hydrogen shipping at $\sim$129{,}500 (68{,}300--202{,}200) averted premature deaths
annually and US\$306\,bn (146--488) saved --- but \textbf{75\% of that is ozone}, only
25\% PM$_{2.5}$. Read as a PM result it is much weaker than the headline suggests.

\item \textbf{Empty clusters.} \textit{Respiratory \& allergic} and \textit{Reproductive
\& developmental} are both unoccupied; respiratory was last empty on 23 August.
\end{itemize}

\clearpage

% ==================================================================== ANALYTICS
\section*{\sffamily\bfseries\color{Deep}Corpus at a glance}
\vspace{-1.5mm}{\color{Amber}\rule{\linewidth}{1.1pt}}\vspace{3mm}

\noindent\includegraphics[width=\linewidth]{\FIGDIR/f1_subtopics.png}
\figcap{Records per subtopic cluster, 19 records over \textbf{eight of ten} clusters.
\textit{Exposure assessment \& modelling} leads with 5, followed by \textit{Sensing} and
\textit{Burden, policy \& mitigation} at 4 each --- so \textbf{13 of 19 records
(68\%) sit on the measurement/modelling/policy side} rather than the clinical side.
\textit{Respiratory \& allergic} and \textit{Reproductive \& developmental} are empty.}

\vspace{2mm}
\noindent\includegraphics[width=\linewidth]{\FIGDIR/f2_design_endpoint.png}
\figcap{Left --- study architecture. \textit{Measurement campaign} 5,
\textit{Experimental / toxicology} 3, \textit{Modelling / inventory} 3,
\textit{Ecological} 2, \textit{Observational --- acute} 2, \textit{Observational ---
cross-sectional} 2, \textit{Observational --- cohort} 1, \textit{Review / synthesis} 1.
The \textit{Experimental} bar is heterogeneous by construction: it pools a murine
instillation study with a wave-channel aerosol-generation experiment and a dark-field
microscopy method paper, because the grouper has no bench-instrumentation bin. Right ---
health endpoint. \textbf{Eight records carry no human endpoint} --- three instrument, three
exposure-construction and two source/bioaerosol --- so 42\% of the corpus is
methodological. No endpoint appears more than once.}

\vspace{2mm}
\noindent\includegraphics[width=\linewidth]{\FIGDIR/f3_metric_geography.png}
\figcap{Left --- particle metric. \textit{PM$_{2.5}$ only} 8,
\textit{PM (unspeciated) / emissions} 6, \textit{PM$_{2.5}$ + PM$_{10}$ jointly} 3, and
one each for \textit{size-resolved / number distribution} (Leibensperger et al., SSA
0.008--20\,$\mu$m) and \textit{composite air-quality index}. Note the recurring
limitation: both speciation-driven records of the day --- Li et al.'s SO$_4^{2-}$-weighted
mixture and Wang et al.'s five-constituent CHARLS analysis --- bin as
\textit{PM$_{2.5}$ only}, because the grouper still has no speciation bin. Right ---
geography: China 3, North America 3, Southeast Asia 3, multi-region 3, Europe 2, East
Asia ex-China 2, and one each from South Asia, Latin America and Oceania. \textbf{Nothing
from Africa}; the last African record was 20 August, eight issues ago.}

\vspace{2mm}
\noindent\includegraphics[width=\linewidth]{\FIGDIR/f4_heatmap.png}
\figcap{Subtopic $\times$ architecture cross-tabulation. Nineteen records over eight
clusters and eight architectures gives a near-diagonal sparsity map: \textbf{four cells
hold 2} --- \textit{Sensing} $\times$ \textit{Measurement campaign},
\textit{Exposure assessment} $\times$ \textit{Measurement campaign},
\textit{Exposure assessment} $\times$ \textit{Modelling / inventory} and
\textit{Burden} $\times$ \textit{Observational --- cross-sectional} --- and the remaining
eleven occupied cells hold 1. Read as sparsity, not structure.}

\vspace{2mm}
\noindent\includegraphics[width=\linewidth]{\FIGDIR/f5_forest.png}
\figcap{Four interval-bearing estimates from two studies. \textbf{The two Li et al.\
estimates are not on a common contrast} --- 1.019 is per IQR of total PM$_{2.5}$ mass,
1.116 is per quartile of the four-component constituent mixture --- so the visual gap
between them is a contrast difference as much as a mixture-versus-mass difference, and
should not be read as an effect ratio. The two Wu et al.\ estimates \emph{are} mutually
comparable in construction (same case-crossover, same Poisson offset structure) but the
PM$_{2.5}$ and NO$_2$ contrasts differ in unit (per 10\,$\mu$g/m$^3$ vs.\ per 10\,ppb);
the NO$_2$ row is carried only as a within-study comparator. The day's other
health results --- the CHARLS sleep-duration coefficients and the Taiwanese mood-disorder
case-crossover --- report on scales that do not belong on a ratio axis and are excluded.}

\vspace{2mm}
\noindent\includegraphics[width=\linewidth]{\FIGDIR/f6_lifecourse.png}
\figcap{Life-course windows, counted only where a record's population is explicitly
anchored to a stage. \textbf{Four records qualify across four of five windows}; the
puberty/adolescence window is empty again. The preconception cell holds the Victorian
women-of-reproductive-age ecological linkage, which is a postcode-level correlation
study and does not resolve any individual exposure window --- it occupies the cell on
population definition alone.}

\clearpage

% ==================================================================== DIGEST
\section*{\sffamily\bfseries\color{Deep}Paper-level digest}
\vspace{-1.5mm}{\color{Amber}\rule{\linewidth}{1.1pt}}\vspace{1mm}

{\fontsize{7.4}{9}\selectfont\color{Slate}Ordered by cluster; the 28 August entry window
only.\par}

\band{Sensing, forecasting \& instrumentation\hfill 4 records}{Amber}

\paperentry{Zhang et al.\ 2026h}
{Contrasting nighttime heterogeneous and daytime photochemical aging drive the optical
evolution of black carbon}
{Atmos Chem Phys}
{Integrated SP2 and single-particle aerosol mass spectrometry over an urban Shenzhen
campaign, mean refractory BC \textbf{1.2\,$\mu$g\,m$^{-3}$}, core mass median diameter
155--170\,nm. The coating-to-core mass ratio rises continuously through the diurnal
cycle by two mechanistically distinct routes --- nighttime heterogeneous coating growth
with enhanced nitrate signal, and daytime photochemical aging with rapid nitrate
accumulation followed by sulfate and oxidised organics. Both converge on similar mass
absorption cross section enhancement at 532\,nm: \textbf{$\sim$0.8\,m$^2$\,g$^{-1}$
overnight, $\sim$1.0\,m$^2$\,g$^{-1}$ by day}. No uncertainty is given on the MAC
enhancements and it is a single city over one campaign, so the magnitudes should not be
transported. Directly relevant to distributed sensing: any optical channel with a fixed
absorption or scattering assumption re-expresses this aging cycle as concentration drift
with a diurnal period.}
{10.5194/acp-26-12197-2026}

\paperentry{Marcon et al.\ 2026}
{Citizen science in air pollution monitoring: methodological approaches, effectiveness
and integration with health research}
{Curr Environ Health Rep}
{Critical review of participatory air-quality monitoring --- low-cost sensors, wearables,
community co-creation --- and its uptake into environmental health research. The review
identifies AI-assisted calibration and distributed sensor networks as the methodological
frontier and argues that hyperlocal citizen data materially improves spatial resolution
and risk communication. Its useful contribution is the ordering of the barriers:
\textbf{representativeness and policy uptake, not sensor accuracy}, are named as the
binding constraints. No quantitative synthesis, no pooled performance figures and no
inclusion criteria are reported in the abstract, so it is a framing document rather than
an evidence base. For a distributed low-cost network this is the argument that siting
design deserves as much engineering attention as the calibration pipeline.}
{10.1007/s40572-026-00559-1}

\paperentry{Wongsasuluk et al.\ 2026}
{PM$_{2.5}$ exposure and respiratory health risk for vulnerable populations in an urban
underpass: a predictive MLR approach, Chiang Mai}
{Environ Sci Pollut Res Int}
{Five DustBoy low-cost stations along the Don Kaew road tunnel, August 2022 to July 2023,
yielding \textbf{385{,}622 records} across PM$_{2.5}$, PM$_{10}$, CO, CO$_2$, NO$_2$,
SO$_2$, O$_3$, TVOC, CH$_2$O and meteorology. A stepwise multiple linear regression
stratified by weekday/weekend and day/night gives a high reported goodness of fit, and a
comparative risk assessment converts predicted concentrations into inhaled dose and
relative risk against Thai and WHO exposure factors. The inference-threatening gap is
that \textbf{no co-location against a reference monitor is reported}: a semi-enclosed
tunnel is a high-humidity, high-loading environment, precisely where nephelometric
low-cost sensors over-read most, and every downstream risk number inherits that bias
uncorrected.}
{10.1007/s11356-026-38171-3}

\paperentry{Xiong et al.\ 2026}
{In situ real-time determination of SO$_2$ photochemical oxidation in nanoscale sea salt
aerosols based on dark-field microscopy}
{Atmos Chem Phys}
{Planar-waveguide dark-field microscopy tracks hygroscopic growth factors of single NaCl
nanoparticles under UV irradiation and SO$_2$, with the Zdanovskii--Stokes--Robinson rule
converting growth to composition and hence to reaction rate. Reported first-order rate
constant \textbf{0.6523\,h$^{-1}$ for 100\,nm NaCl}, with a \textbf{non-monotonic size
dependence} --- rising from 50 to 200\,nm, falling above 200\,nm --- attributed to
competition between surface curvature and specific surface area, and cross-checked
against transition-state theory and a double-film mass-transfer treatment. Single
salt system under controlled UV; the extrapolation to ambient mixed-composition marine
aerosol is not demonstrated. Methodologically this is the interesting part: an optical
single-particle technique that yields kinetics without collection, at sizes below what
most PM instrumentation resolves at all.}
{10.5194/acp-26-12183-2026}

\clearpage

\band{Exposure assessment \& modelling\hfill 5 records}{Amber}

\paperentry{Cui et al.\ 2026}
{Daily 1\,km satellite-based PM$_{2.5}$ exposure estimates in Bangladesh, 2020--2025}
{Environ Pollut}
{Two-stage machine learning over one of the most monitoring-sparse high-burden regions in
the world. Stage one reconstructs MAIAC aerosol optical depth, of which
\textbf{more than 60\% of potential retrievals are missing}, with LightGBM; stage two
feeds the gap-filled AOD plus meteorological, composition and geospatial predictors into
an attention-based TabNet to estimate ground PM$_{2.5}$. Independent test performance
\textbf{$R^2 = 0.95$, RMSE 17.41\,$\mu$g/m$^3$, MAE 9.38}, against a national annual mean
of \textbf{68--75\,$\mu$g/m$^3$} and winter peaks above 130. The limitation is structural:
with almost no ground reference network, the held-out set is small and drawn from the
same sparse monitors used in training, so $R^2 = 0.95$ describes interpolation skill more
than transferability, and the DHS linkage is division-level and ecological.}
{10.1016/j.envpol.2026.129051}

\paperentry{Hopke et al.\ 2026}
{A PM$_{2.5}$ magnitude-rank index for summarizing annual concentration-frequency
structure}
{Environ Monit Assess}
{Introduces the PM$_{2.5}$ Magnitude--Rank Index --- the largest integer $k$ such that at
least $k$ days in a year have daily mean PM$_{2.5}$ $\geq k$\,$\mu$g/m$^3$ --- as a
threshold-free joint summary of magnitude and frequency. Applied to \textbf{918
qualifying MSA-years across 96 US metropolitan areas, 2013--2022}, PMRI ranges 6--38
(median 16) and correlates with the annual mean at Spearman $\rho = 0.83$, the annual
98th percentile at 0.95, and days $\geq15$\,$\mu$g/m$^3$ at 0.96. That last figure is the
honest problem with the metric: at $\rho = 0.96$ against an existing exceedance count it
adds very little information, so its case rests on being policy-threshold-free and
intuitively readable rather than on discriminating power. No health-outcome validation is
attempted.}
{10.1007/s10661-026-15858-1}

\paperentry{Reyna-Bello et al.\ 2026}
{Seasonal variability of PM$_{2.5}$-bound metals and meteorology in an urban-industrial
area of northern Mexico}
{Environ Monit Assess}
{MiniVol low-volume sampling across four seasons in Saltillo during 2023, with ten metals
(Al, Cd, Cu, Fe, K, Mn, Na, Ni, Pb, Zn) by flame atomic absorption and PCA against
temperature, humidity, wind, precipitation and pressure. Mass is strongly cold-season
weighted --- \textbf{152.9$\pm$23.6\,$\mu$g/m$^3$ in autumn and 115.3$\pm$37.2 in winter
against 40.7$\pm$6.8 in spring} --- consistent with suppressed dispersion, while several
metals invert that pattern and peak in spring and summer. The reported metal
concentrations are implausibly high for PM$_{2}$$_{.5}$-bound species (Cd at
3.13\,$\mu$g/m$^3$ is orders of magnitude above typical urban values), which points to a
units or digestion-blank problem; treat the metal magnitudes as unverified and the
seasonal contrast as the usable result.}
{10.1007/s10661-026-15856-3}

\paperentry{Leibensperger et al.\ 2026}
{Sea surface warming suppresses primary sea spray aerosol number production}
{Atmos Chem Phys}
{Controlled wave-breaking experiments in the Scripps Ocean-Atmosphere Research Simulator,
which reproduces wind, waves and sea surface temperature together. Raising SST from 2 to
23\,$^\circ$C \textbf{suppressed subsurface bubble concentration by 1.5$\times$, sea
spray aerosol number by 3$\times$ and SSA number flux by 3$\times$}, with a
non-monotonic SST dependence in the supermicron fraction. This is a single simulator with
one seawater matrix and no biological aging, so absolute fluxes are not transferable ---
the value is the sign and rough magnitude of the SST dependence, which the field has not
agreed on. Relevance here is indirect but real: marine background number concentration is
the baseline against which coastal low-cost network anomalies are read.}
{10.5194/acp-26-12211-2026}

\paperentry{Agpaoa et al.\ 2026}
{Occurrence, temporal variation and meteorological drivers of airborne fungal spores in
Laoag City, northern Philippines}
{Environ Monit Assess}
{First aeromycological record for a monsoon-influenced coastal site in northern Luzon:
216 sampling days, March to November 2022, cumulative index \textbf{190{,}096
spores\,m$^{-3}$}, dominated by \textit{Cladosporium} and Fusaroid-like types. Wet-season
sampling carries most of the cumulative load but the highest daily concentrations fall in
the dry season, and diurnal peaks sit at 00:00--06:00. Positive associations with relative
humidity and rainfall, negative with temperature, and no significant wind association.
Single site, single year, morphology-based identification only, and no PM co-measurement
--- so the bioaerosol fraction cannot be placed within a mass-based PM budget. Included
because overnight spore peaks are a documented confounder for optical particle counters
in humid climates.}
{10.1007/s10661-026-15814-z}

\clearpage

\band{Burden, policy \& mitigation\hfill 4 records}{Coral}

\paperentry{Zhang et al.\ 2026g}
{Health and economic impacts of improved air quality from transitioning global shipping
energy to hydrogen}
{Nat Commun}
{Integrated emission--air quality--health framework over feasible hydrogen-transition
scenarios for global shipping. Full adoption averts approximately \textbf{129{,}500
(95\% CI 68{,}300--202{,}200) premature deaths annually}, valued at \textbf{US\$306\,bn
(146--488)}, with India, Japan and the United States taking the largest economic
benefit. The number that matters for a PM watch is the split: \textbf{75\% of the averted
mortality is ozone, only 25\% PM$_{2.5}$}, with the O$_3$ gains concentrated in Southern
Asia and the PM$_{2.5}$ gains in Eastern Asia. The confidence interval spans a factor of
three and propagates concentration--response uncertainty only, not scenario or
inventory uncertainty, so the true spread is wider than stated.}
{10.1038/s41467-026-76134-x}

\paperentry{Nilsson et al.\ 2026}
{Urban environment, health and equity: symptoms and annoyances related to air pollution,
traffic noise and greenness}
{Environ Epidemiol}
{\textbf{46{,}282 adults aged 18--84} in Stockholm County responding to the 2023 Swedish
National Environmental Health Survey, with geocoded residential PM$_{2.5}$, PM$_{10}$,
NO$_2$, road/rail/aircraft noise and 300\,m NDVI, dichotomised and stratified by
socioeconomic status and building type. Apartment residents are more exposed to pollution
and noise and have less greenness; SES-stratified patterns are exposure-specific with no
consistent overall gradient. The strong associations are with \emph{perception} --- NO$_2$
with poor perceived outdoor air quality (\textbf{OR 3.22; 95\% CI 2.75--3.79}), road noise
with high annoyance (OR 4.59; 4.02--5.26). Both exposure and outcome are cross-sectional
and the outcome is self-reported, so these are perception--exposure concordance
estimates, not health effects, and the odds ratios should not be read as such.}
{10.1097/EE9.0000000000000516}

\paperentry{Sirikul et al.\ 2026}
{Knowledge, attitudes and practices of caregivers regarding PM$_{2.5}$ exposure among
children in northern Thailand}
{J Clim Chang Health}
{Cross-sectional KAP survey of caregivers recruited at Chiang Mai Child Development
Centres. Knowledge is moderate (10.44$\pm$3.80 of 22), attitudes positive
(30.74$\pm$3.57 of 40), practices inconsistent (17.33$\pm$5.22 of 32). The specific gaps
are the useful part: \textbf{97.0\% recognise the health risk to children, but only
41.8\% know surgical masks are inadequate for PM$_{2.5}$ and only 9.7\% can correctly
interpret an AQI value}; air purifier ownership is 15.3\% and monitoring-app use 28.2\%,
both with a clear education and income gradient. The abstract does not report the sample
size, response rate or instrument validation, and centre-based recruitment excludes
non-attending households, so absolute prevalences are not generalisable. The AQI-literacy
figure is the one that should worry anyone designing a public-facing sensor dashboard.}
{10.1016/j.joclim.2026.100734}

\paperentry{Sharma et al.\ 2026b}
{Geospatial linkage of environmental exposures, socioeconomic profile and health outcomes
among women of reproductive age, Victoria}
{Health Promot J Austr}
{Postcode-level linkage of 2021 Australian Census sociodemographic and health variables to
PM$_{2.5}$, NO$_2$, CO, SO$_2$, formaldehyde, ozone, land surface temperature and
greenness, with Moran's I, Wagstaff concentration indices and geospatial regression.
Non-communicable disease burden in women of reproductive age is strongly socioeconomically
patterned (\textbf{323.0 vs.\ 227.3 per 1000} in most- versus least-disadvantaged
postcodes; steepest for diabetes, Wagstaff $-0.16$), \textbf{but environmental exposures
show only weak associations with outcomes} despite clear spatial clustering. Inner
Melbourne has lower NCD burden despite elevated NO$_2$ and PM$_{2.5}$. This is
ecological, cross-sectional and confounded by the same SES gradient it measures; the
null environmental finding is uninformative rather than reassuring.}
{10.1002/hpja.70241}

\clearpage

\band{Neuro / mental health\hfill 2 records}{Deep}

\paperentry{Wang et al.\ 2026e}
{Speciated fine particulate matter and self-reported sleep time in middle-aged and older
Chinese adults}
{Sleep Med}
{\textbf{19{,}505 CHARLS participants aged 45+}, 2011--2018, with 6-month and 2-year
PM$_{2.5}$ and five constituents --- BC, organic matter, SO$_4^{2-}$, NO$_3^-$,
NH$_4^+$ --- from the Tracking Air Pollution in China product, against total sleep time
(nocturnal plus napping) via linear mixed models and quantile g-computation. The
secondary inorganic ions NO$_3^-$, SO$_4^{2-}$ and NH$_4^+$ show the most consistent
associations, and the direction is the surprise: highest versus lowest quartile of
intermediate-term exposure \textbf{increases} total sleep time (PM$_{2.5}$ +0.16\,h and
comparable for the constituents). Self-reported duration, no actigraphy, and an
increase in sleep duration is not obviously a harm --- it may index fatigue, reduced
outdoor activity on high-pollution days, or reverse behavioural adaptation. Mechanism is
not addressed and the finding should be treated as hypothesis-generating.}
{10.1016/j.sleep.2026.109229}

\paperentry{Chen et al.\ 2026i}
{Gender and season disparities in the influence of short-term PM$_{2.5}$ exposure on
emergency room visits for mood disorders in Taiwan}
{J Occup Environ Med}
{Time-stratified case-crossover with distributed lag non-linear models over
\textbf{2{,}172 emergency visits} for bipolar or depressive disorder in the Chang Gung
Research Database, 2010--2020, linked to Ministry of Environment monitoring. Associations
appear \textbf{exclusively in the hot season and predominantly in women}, with delayed
effects at lags 3--5 and cumulative effects over lags 0--1 and 0--3 to 0--5. No effect
estimates, confidence intervals or exposure contrasts appear in the abstract, which makes
the result unquantifiable as reported. With 2{,}172 events split by sex and season the
subgroup cells are small, and a hot-season-only, women-only finding from a design with no
pre-registered subgroup plan is the classic shape of a multiplicity artefact. Carried at
tier B on design quality, not on the strength of the reported signal.}
{10.1097/JOM.0000000000003848}

\clearpage

\band{Cardiovascular \& metabolic\hfill 1 record}{Teal}

\paperentry{Li et al.\ 2026e}
{Sulfate-dominated PM$_{2.5}$ constituents are associated with increased acute myocardial
infarction mortality: a US nationwide difference-in-differences study}
{Environ Sci Technol Lett}
{County-level US AMI mortality against PM$_{2.5}$ and its EC, OC, SO$_4^{2-}$ and
NO$_3^-$ constituents, analysed with a causal mixture method inside a
difference-in-differences design. Total mass: \textbf{RR 1.019 (95\% CI 1.013--1.025) per
IQR}. Constituent mixture: \textbf{RR 1.116 (1.101--1.130) per one-quartile increase},
with \textbf{SO$_4^{2-}$ carrying weight 0.762}. Regionally, an eastern county sees
\textbf{1.545 (0.656--2.437) excess AMI deaths per 100{,}000 per year} relative to the
west per quartile of mixture, SO$_4^{2-}$-driven in the east (0.707) and OC-driven in the
west (0.485). Exposure is county-level and the outcome is mortality rather than incidence,
so this is ecological in structure --- the DiD design controls time-invariant county
confounding but not time-varying composition-correlated confounders such as regional
industrial decline.}
{10.1021/acs.estlett.6c00288}

\band{Mechanistic toxicology\hfill 1 record}{Violet}

\paperentry{Van et al.\ 2026}
{MASLD exacerbates chronic low-dose PM$_{2.5}$-induced lung injury, inflammation and
fibrosis}
{In Vivo}
{C57BL/6J mice on a Western or normal diet for 28 weeks with concurrent intratracheal
diesel particulate matter at \textbf{0.5\,mg/kg three times weekly} (stated as
environmentally relevant to $\sim$50\,$\mu$g/m$^3$), assessed by histology, Masson's
trichrome, Luminex multiplex on bronchoalveolar lavage and lung-tissue ELISA/qPCR. The
Western diet induced the MASLD phenotype; PM$_{2.5}$ did not worsen the hepatic
phenotype, but \textbf{the combination markedly exacerbated pulmonary injury and fibrosis
relative to PM$_{2.5}$ alone}. The dose-equivalence claim is the weak point: bolus
intratracheal instillation of diesel PM delivers a deposition profile nothing like
chronic inhalation at 50\,$\mu$g/m$^3$, so the ``environmentally relevant'' framing should
not be taken at face value. The interaction direction is the transferable result.}
{10.21873/invivo.14428}

\clearpage

\band{Occupational \& indoor\hfill 1 record}{Clay}

\paperentry{Ospalek et al.\ 2026}
{Air quality assessment and potential inhalation exposure across equestrian facilities}
{Atmosphere}
{Nine measurements at 1.5\,m in a stable, an indoor riding arena and an outdoor riding
arena at a Slovak university equestrian centre, covering PM$_1$, PM$_{2.5}$, PM$_{10}$,
respirable and total dust plus gaseous pollutants, with four exposure scenarios combining
18--20\,h stabled with 4--6\,h ridden. The \textbf{outdoor arena carried the highest
concentrations of every PM fraction and the indoor arena the lowest} ($p<0.05$) --- the
opposite of the expected enclosure effect, and the paper's most interesting result;
ammonia was higher in the stable, as expected. Summer measurements only, one facility,
nine measurements total, and the exposure scenarios are constructed rather than observed,
so the inhaled-dose estimates are illustrative. Relevant here as a template for
microenvironment-resolved dose reconstruction from sparse fixed-point sampling.}
{10.3390/atmos17090840}

\band{Other clinical endpoints\hfill 1 record}{Slate}

\paperentry{Wu et al.\ 2026d}
{Short-term environmental exposures and emergency department imaging utilization}
{J Am Coll Radiol}
{Time-stratified case-crossover over \textbf{1{,}946{,}465 imaging studies} at five
Toronto academic hospitals, 2013--2023, with conditional Poisson models against
PM$_{2.5}$, NO$_2$, temperature, wind, rainfall and snowfall, modelling both absolute
volume and per-visit rate (log daily visits as offset). Higher PM$_{2.5}$ associates with
higher imaging volume, \textbf{IRR 1.020 (95\% CI 1.017--1.024) per 10\,$\mu$g/m$^3$},
alongside NO$_2$ 1.017 (1.014--1.020) per 10\,ppb and temperature 1.040 per 10\,$^\circ$C;
wind, rain and snow are protective, snowfall most strongly (0.959 per 10\,mm). The
volume effect is almost certainly mediated by visit counts rather than by per-visit
imaging intensity, and the exposures are collinear with each other and with season ---
so this is a health-system demand signal, not evidence of a PM-specific clinical
mechanism. Its value is operational: PM$_{2.5}$ forecasts carry usable information for ED
resource planning.}
{10.1016/j.jacr.2026.08.014}

\clearpage

% ==================================================================== TABLE
\section*{\sffamily\bfseries\color{Deep}Corpus register}
\vspace{-1.5mm}{\color{Amber}\rule{\linewidth}{1.1pt}}\vspace{2.5mm}

{\fontsize{6.7}{8.4}\selectfont
\setlength{\tabcolsep}{2.4pt}
\begin{longtable}{@{}>{\RaggedRight}p{25mm}>{\RaggedRight}p{41mm}>{\RaggedRight}p{27mm}%
>{\RaggedRight}p{23mm}>{\RaggedRight}p{18mm}>{\RaggedRight}p{15mm}@{}}
\toprule
\textsf{\textbf{First author}} & \textsf{\textbf{Focus}} & \textsf{\textbf{Journal}} &
\textsf{\textbf{Design}} & \textsf{\textbf{Metric}} & \textsf{\textbf{Region}}\\
\midrule
\endfirsthead
\toprule
\textsf{\textbf{First author}} & \textsf{\textbf{Focus}} & \textsf{\textbf{Journal}} &
\textsf{\textbf{Design}} & \textsf{\textbf{Metric}} & \textsf{\textbf{Region}}\\
\midrule
\endhead
\multicolumn{6}{@{}l@{}}{\cellcolor{Mist}\textsf{\textbf{\textcolor{Deep}{Neuro / mental health}}}}\\[0.6mm]
\href{https://doi.org/10.1001/jamanetworkopen.2026.25015}{Dales} & Wildfire-sourced PM$_{2.5}$ and ED visits for migraine and other primary headache & \textit{JAMA Netw Open} & Case-crossover & PM$_{2.5}$ (wildfire) & Canada\\
\href{https://doi.org/10.1093/ageing/afag215}{Glover} & Umbrella meta-review of urban environmental exposures and cognitive health & \textit{Age Ageing} & Umbrella review / meta-analysis & PM$_{2.5}$ & Global\\
\href{https://doi.org/10.1016/j.nsa.2026.107019}{Goossens} & In utero air pollution exposure and pre-/postnatal brain development & \textit{Neurosci Appl} & Narrative review & PM$_{2.5}$ / PM$_{10}$ / PM$_{\mathrm{coarse}}$ & Global\\
\href{https://doi.org/10.1016/j.isci.2026.116617}{Ma} & PM$_{2.5}$-induced neurotoxicity and depression: exposure risks and mechanisms & \textit{iScience} & Narrative review & PM$_{2.5}$ + components & Global\\
\href{https://doi.org/10.3760/cma.j.cn112150-20250804-00752}{Pan} & Maternal PM exposure and offspring neurodevelopment: effects and mechanisms & \textit{Zhonghua Yu Fang Yi Xue Za Zhi} & Narrative review & PM (bulk) & China\\
\addlinespace[0.9mm]
\multicolumn{6}{@{}l@{}}{\cellcolor{Mist}\textsf{\textbf{\textcolor{Teal}{Cardiovascular \& metabolic}}}}\\[0.6mm]
\href{https://doi.org/10.3760/cma.j.cn112338-20251120-00843}{Huang} & PM$_{2.5}$ and dyslipidaemia in rural Chinese children; dietary effect modification & \textit{Zhonghua Liu Xing Bing Xue Za Zhi} & Cross-sectional (multistage) & PM$_{2.5}$ & China\\
\href{https://doi.org/10.32394/pe/218966}{Kazmierska} & Air pollution and cardiovascular disease: review of current data & \textit{Przegl Epidemiol} & Narrative review & PM$_{2.5}$ / PM$_{10}$ & Global\\
\href{https://doi.org/10.1007/s00125-026-06814-2}{Xi} & Air pollution, built environment and temperature vs. prediabetes phenotypes & \textit{Diabetologia} & Prospective cohort & PM$_{2.5}$ & Germany\\
\href{https://doi.org/10.1161/JAHA.125.046917}{Yang} & Short-term PM$_{2.5}$ and hospitalisation/ED use in atrial fibrillation & \textit{J Am Heart Assoc} & Case-crossover & PM$_{2.5}$ (wildfire-influenced) & USA\\
\addlinespace[0.9mm]
\multicolumn{6}{@{}l@{}}{\cellcolor{Mist}\textsf{\textbf{\textcolor{Sky}{Reproductive \& developmental}}}}\\[0.6mm]
\href{https://doi.org/10.1016/j.jes.2025.10.043}{Chen} & Prenatal air pollution, fetal growth trajectory phenotypes and child neurodevelopment & \textit{J Environ Sci (China)} & Prospective cohort & PM$_{2.5}$ / PM$_{10}$ & China\\
\href{https://doi.org/10.1016/j.ecoenv.2026.120566}{Huang} & Joint prenatal bisphenol + air-pollutant mixture exposure and preschool IQ & \textit{Ecotoxicol Environ Saf} & Prospective cohort (mixtures) & PM$_{2.5}$ / PM$_{10}$ & China\\
\href{https://doi.org/10.1007/s40572-026-00552-8}{Peterson} & Ambient air pollution, PAHs and pubertal development: critical review & \textit{Curr Environ Health Rep} & Critical review & PM$_{2.5}$ components & Global\\
\href{https://doi.org/10.1016/j.envres.2026.125313}{Xie} & Long-term PM$_{2.5}$ and timing of pubertal development in boys and girls & \textit{Environ Res} & Prospective cohort & PM$_{2.5}$ & China\\
\addlinespace[0.9mm]
\multicolumn{6}{@{}l@{}}{\cellcolor{Mist}\textsf{\textbf{\textcolor{Sage}{Respiratory \& allergic}}}}\\[0.6mm]
\href{https://doi.org/10.1007/s00484-026-03289-z}{Bayar Muluk} & Climate change and otolaryngology: airway disease impacts and emission reduction & \textit{Int J Biometeorol} & Narrative review & PM$_{2.5}$ & Global\\
\href{https://doi.org/10.4103/mgr.MEDGASRES-D-25-00164}{Liao} & Short-term air pollution, respiratory symptoms and lung function in asthma & \textit{Med Gas Res} & Panel study & PM$_{2.5}$ / PM$_{10}$ & China\\
\href{https://doi.org/10.3390/epidemiologia7040097}{Sanchez de Prada} & Ambient PM and asthma-coded admissions in children 0-3 y, Spain & \textit{Epidemiologia} & Ecological panel & PM$_{2.5}$ / PM$_{10}$ & Spain\\
\href{https://doi.org/10.1016/j.jhazmat.2026.143076}{Yang} & Allergenic vs. toxic risk of airborne seed hairs compared with ambient PM$_{2.5}$ & \textit{J Hazard Mater} & Field measurement & PM$_{2.5}$ + bioaerosol & China\\
\addlinespace[0.9mm]
\multicolumn{6}{@{}l@{}}{\cellcolor{Mist}\textsf{\textbf{\textcolor{Violet}{Mechanistic toxicology}}}}\\[0.6mm]
\href{https://doi.org/10.1016/j.jes.2025.12.074}{Ding} & PM$_{2.5}$ promotes vascular calcification via miR-27b-5p/HES1 signalling & \textit{J Environ Sci (China)} & In vitro + organoid + murine & PM$_{2.5}$ & China\\
\href{https://doi.org/10.1007/s00360-026-01698-5}{Maraschi} & Osmoregulatory disruption and metal bioaccumulation from settleable atmospheric PM & \textit{J Comp Physiol B} & Controlled exposure (ecotox) & Settleable APM & Brazil\\
\href{https://doi.org/10.1016/j.jes.2025.11.018}{Wu} & Macrophage NLRP3 inflammasome in real-ambient PM-induced hepatic glucose dysregulation & \textit{J Environ Sci (China)} & Real-ambient murine + co-culture & Real-ambient PM & China\\
\addlinespace[0.9mm]
\multicolumn{6}{@{}l@{}}{\cellcolor{Mist}\textsf{\textbf{\textcolor{Clay}{Occupational \& indoor}}}}\\[0.6mm]
\href{https://doi.org/10.3389/fmedt.2026.1796828}{Chen} & Nested suction instrument to cut intraoperative surgical-smoke PM$_{2.5}$ & \textit{Front Med Technol} & Bench / device experiment & PM$_{2.5}$ (surgical smoke) & China\\
\href{https://doi.org/10.1093/annweh/wxag055}{Cole} & Determinants of respirable crystalline silica exposure in Australian tunnelling & \textit{Ann Work Expo Health} & Observational exposure study & Respirable dust / RCS & Australia\\
\href{https://doi.org/10.1016/j.jhazmat.2026.143081}{Ding} & Stage-resolved hair Zn-isotope dynamics as a firefighting combustion-exposure tracer & \textit{J Hazard Mater} & Repeated-measures biomarker & Combustion-derived PM & China\\
\href{https://doi.org/10.3390/jox16040119}{Ramsden} & Xenobiotic hazards in aircraft cabin air & \textit{J Xenobiot} & Review + risk economics & Ultrafine particles & Global\\
\addlinespace[0.9mm]
\multicolumn{6}{@{}l@{}}{\cellcolor{Mist}\textsf{\textbf{\textcolor{Amber}{Sensing, forecasting \& instrumentation}}}}\\[0.6mm]
\href{https://doi.org/10.1016/j.jenvman.2026.130515}{Shi} & Adaptive spatiotemporal graph transformer for multi-site multi-step PM$_{2.5}$ forecasting & \textit{J Environ Manage} & Deep-learning model & PM$_{2.5}$ & China / India\\
\href{https://doi.org/10.3390/bios16070380}{Tsuchiya} & Seasonal and indoor-outdoor characteristics of fluorescent aerosol particles in offices & \textit{Biosensors} & Real-time sensor campaign & Bioaerosol / coarse & Japan\\
\addlinespace[0.9mm]
\multicolumn{6}{@{}l@{}}{\cellcolor{Mist}\textsf{\textbf{\textcolor{Amber}{Exposure assessment \& modelling}}}}\\[0.6mm]
\href{https://doi.org/10.1029/2025GH001621}{Solorzano-Araque} & Multiscale exposure modelling: LOTOS-EUROS + mobility surveys + particle toxicity & \textit{GeoHealth} & CTM + mobility integration & PM$_{2.5}$ / PM$_{10}$ & Colombia\\
\addlinespace[0.9mm]
\multicolumn{6}{@{}l@{}}{\cellcolor{Mist}\textsf{\textbf{\textcolor{Coral}{Burden, policy \& mitigation}}}}\\[0.6mm]
\href{https://doi.org/10.1002/hsr2.72886}{Karamzad} & Ischaemic heart disease burden in MENA 1990-2021; PM attributable fraction 28.2\% & \textit{Health Sci Rep} & GBD secondary analysis & Ambient PM & MENA\\
\href{https://doi.org/10.1371/journal.pone.0347989}{Rees} & HEPA filtration to reduce respiratory infection in care homes & \textit{PLoS One} & RCT process evaluation & Filtered PM & UK\\
\href{https://doi.org/10.1016/j.jes.2026.01.009}{Shan} & Policy-driven emission reductions and post-pandemic rebound, high-resolution inventories & \textit{J Environ Sci (China)} & Emission inventory & PM emissions & China\\
\addlinespace[0.9mm]
\multicolumn{6}{@{}l@{}}{\cellcolor{Mist}\textsf{\textbf{\textcolor{Slate}{Other clinical endpoints}}}}\\[0.6mm]
\href{https://doi.org/10.3390/dj14070440}{Lupita} & PM$_{2.5}$, preventive-care capacity and caries/edentulism burden across the EU & \textit{Dent J} & Ecological panel & PM$_{2.5}$ & EU-27\\
\href{https://doi.org/10.1016/j.jdin.2026.06.008}{Park} & Long-term PM exposure and risk of androgenetic alopecia & \textit{JAAD Int} & Nationwide cohort & PM$_{2.5}$ / PM$_{10}$ & South Korea\\
\href{https://doi.org/10.3171/2026.3.PEDS2615}{Sun} & Environmental exposure and area deprivation vs. paediatric hydrocephalus outcomes & \textit{J Neurosurg Pediatr} & Registry cohort + geospatial & PM$_{2.5}$ (satellite-derived) & USA\\
\addlinespace[0.9mm]
\bottomrule
\end{longtable}}


\vspace{4mm}
\begin{keybox}
{\sffamily\bfseries\footnotesize\color{Deep}Provenance and method}\\[1.2mm]
{\fontsize{7.2}{8.8}\selectfont
Entry window \textbf{2026-08-28 to 2026-08-28}, a single day of deposits.
\textbf{Sources.} PubMed E-utilities on \texttt{[Date - Entry]}, run as two separate
queries (health axis 17, instrumentation axis 2, one record shared). Europe PMC
\texttt{CREATION\_DATE} returned 0 for this date. Crossref by-ISSN across the tracked
non-PubMed-indexed titles returned 57 deposits. OpenAlex returned HTTP 429 and remains
broken for entry-date windowing; arXiv contributed nothing dated in window.
\textbf{Screening.} 76 raw records, \textbf{73 unique} after PMID/DOI collapse (one
PubMed record returned on both axes, one Crossref deposit duplicated within the feed, and
the Bangladesh paper arrived through both PubMed and Crossref), \textbf{0 already in}
\texttt{seen.json}, \textbf{19 retained and 54 rejected} with reason logged to
\texttt{state/rejected.jsonl}. Rejections break down as 43 out of scope (non-particle
atmospheric chemistry, water and soil contamination, NDT and metrology deposits from
\textit{Meas Sci Technol}), 8 in scope by title but with \textbf{no retrievable
abstract} --- carried to the retry list rather than summarised --- and 3 editorial front
matter (two retraction notices, one corrigendum).
\textbf{Effect estimates.} Four estimates from two studies; the figure caption states why
none of the four pairs is fully comparable.
\textbf{Pipeline changes made this run.} Seventeen design labels and \textit{Slovakia}
were added to \texttt{plots.py} in the same edit that wrote them onto records, per the
22 August guard.
\textbf{Labels.} Subtopic, design, particle-metric and geography are single-label by
dominant contribution --- a judgement, not a controlled vocabulary. All 19 DOIs passed
\texttt{check\_dois.py} (0 fail, 0 warn) before build.\par}
\end{keybox}

\end{document}
